% ============================================================================
%  TAVS-ESP : formal specification
%
%  These equations describe the system AS IMPLEMENTED, which differs from the
%  original manuscript in four places. Each divergence is marked with a
%  \changed{} note giving the reason, so the corresponding manuscript sections
%  can be revised rather than silently contradicted.
%
%  Standalone: pdflatex tavs_equations.tex
%  Or drop into a thesis with % ============================================================================
%  TAVS-ESP : formal specification
%
%  These equations describe the system AS IMPLEMENTED, which differs from the
%  original manuscript in four places. Each divergence is marked with a
%  \changed{} note giving the reason, so the corresponding manuscript sections
%  can be revised rather than silently contradicted.
%
%  Standalone: pdflatex tavs_equations.tex
%  Or drop into a thesis with % ============================================================================
%  TAVS-ESP : formal specification
%
%  These equations describe the system AS IMPLEMENTED, which differs from the
%  original manuscript in four places. Each divergence is marked with a
%  \changed{} note giving the reason, so the corresponding manuscript sections
%  can be revised rather than silently contradicted.
%
%  Standalone: pdflatex tavs_equations.tex
%  Or drop into a thesis with % ============================================================================
%  TAVS-ESP : formal specification
%
%  These equations describe the system AS IMPLEMENTED, which differs from the
%  original manuscript in four places. Each divergence is marked with a
%  \changed{} note giving the reason, so the corresponding manuscript sections
%  can be revised rather than silently contradicted.
%
%  Standalone: pdflatex tavs_equations.tex
%  Or drop into a thesis with \input{tavs_equations} (then delete the preamble
%  and \begin{document}/\end{document}).
% ============================================================================
\documentclass[11pt]{article}

\usepackage[margin=1in]{geometry}
\usepackage{amsmath,amssymb,amsthm}
\DeclareMathOperator*{\argmin}{arg\,min}  % not provided by amsmath
\usepackage{booktabs}
\usepackage{xcolor}
\usepackage[colorlinks=true,linkcolor=black,citecolor=black]{hyperref}

% Divergences from the original formulation are flagged, not hidden.
\newcommand{\changed}[1]{%
  \par\vspace{2pt}\noindent
  \textcolor{red!65!black}{\small\textbf{Differs from manuscript:} #1}\par\vspace{2pt}}

% ---- notation shorthands -----------------------------------------------------
\newcommand{\round}{r}
\newcommand{\Vset}{\mathcal{V}}          % verified
\newcommand{\Pset}{\mathcal{P}}          % promoted
\newcommand{\Dset}{\mathcal{D}}          % decoys
\newcommand{\Lset}{\mathcal{L}}          % inliers
\newcommand{\Oset}{\mathcal{O}}          % outliers
\newcommand{\Teff}{\widetilde{T}}        % effective (ramp-capped) trust
\newcommand{\gproj}{\hat{g}}             % projected update
\newcommand{\sigsq}{\hat{\sigma}}
\newcommand{\1}[1]{\mathbf{1}\!\left[#1\right]}  % amssymb \mathbb covers uppercase only

\title{\textbf{TAVS--ESP: Formal Specification}\\[2pt]
       \large Trust-Adaptive Verification Scheduling with Ephemeral Structured Projection}
\author{}
\date{}

\begin{document}
\maketitle
\vspace{-3em}

% =============================================================================
\section{Notation}

\begin{center}
\begin{tabular}{@{}ll@{\qquad}ll@{}}
\toprule
$\round$ & round index & $i$ & client index \\
$m$ & parameter block index & $g_i$ & parameters submitted by client $i$ \\
$n_i$ & size of client $i$'s dataset & $d_m$ & dimension of block $m$ \\
$T_i(\round)$ & trust score & $\Teff_i(\round)$ & effective (ramp-capped) trust \\
$c_i$ & consecutive clean verifications & $a_i$ & appearances since last verified \\
$\ell_i$ & round of last verification & $r_0^i$ & round client $i$ joined \\
$\Vset(\round)$ & verified set & $\Pset(\round)$ & promoted set \\
$\Dset(\round)\subseteq\Vset(\round)$ & decoys & $\Lset(\round)$ & inliers \\
\bottomrule
\end{tabular}
\end{center}

\medskip
\noindent
Parameters: $\theta_{\mathrm{low}},\theta_{\mathrm{high}}$ (tier thresholds),
$\alpha$ (trust EMA), $\tau_{\mathrm{ramp}}$, $k_{\mathrm{trust}}$,
$p_{\mathrm{decoy}}$, $c_\lambda$, $\gamma_{\mathrm{budget}}$,
$s^{\mathrm{app}}_{\max}$, $s^{\mathrm{rnd}}_{\max}$,
$\tau_z$, $\alpha_\sigma$, $\kappa$ (clip factor), $k$ (projection budget).

% =============================================================================
\section{Layer 1 --- Trust-Adaptive Verification Scheduling}

\subsection{Effective trust}

A newly joined client cannot accrue influence quickly, regardless of behaviour
(Sybil resistance):
\begin{align}
T^{\max}_i(\round) &= 1 - \exp\!\left(-\frac{\round - r_0^i}{\tau_{\mathrm{ramp}}}\right),
\label{eq:ramp}\\[2pt]
\Teff_i(\round)    &= \min\!\bigl(T_i(\round),\, T^{\max}_i(\round)\bigr).
\label{eq:efftrust}
\end{align}

\subsection{Tier assignment}

\begin{equation}
\mathrm{tier}_i(\round)=
\begin{cases}
1, & \Teff_i(\round) < \theta_{\mathrm{low}},\\[2pt]
3, & \Teff_i(\round) \ge \theta_{\mathrm{high}} \ \wedge\  c_i \ge k_{\mathrm{trust}},\\[2pt]
2, & \text{otherwise.}
\end{cases}
\label{eq:tier}
\end{equation}

Tier~3 requires \emph{both} conditions: a trust threshold alone would let a
single favourable round grant promotion.

\subsection{Staleness}

Trust states how well a client behaved when last observed; staleness states
whether that observation is still current. The two are tracked separately.
\begin{equation}
\mathrm{stale}_i(\round)=
\1{a_i \ge s^{\mathrm{app}}_{\max}}
\ \vee\
\1{\round - \ell_i \ge s^{\mathrm{rnd}}_{\max}}
\ \vee\
\1{\text{$i$ never verified}}.
\label{eq:stale}
\end{equation}

Both caps are required: they measure different quantities. At $40\%$
participation, $s^{\mathrm{app}}_{\max}=4$ appearances corresponds to roughly
ten wall-clock rounds, so neither bound implies the other. A client that has
never been verified is stale by definition and cannot be promoted on no
evidence.

\subsection{Assignment rule}

\begin{equation}
i \longmapsto
\begin{cases}
\Vset, & \mathrm{stale}_i(\round) \quad \text{(hard expiry, outranks tier)},\\[2pt]
\Vset, & \mathrm{tier}_i = 1,\\[2pt]
\Vset \ \text{and} \ \Dset, & \mathrm{tier}_i = 3 \ \wedge\ \mathrm{CSPRNG}(\mathrm{key},i,\round) < p_{\mathrm{decoy}},\\[2pt]
\Pset, & \mathrm{tier}_i = 3 \ \wedge\ \text{otherwise},\\[2pt]
\Pset, & \mathrm{tier}_i = 2 \quad \text{(tentative; see \eqref{eq:budget})}.
\end{cases}
\label{eq:assign}
\end{equation}

Decoys are verified server-side while the client is told it was promoted, so no
client can determine whether it is being observed. The expiry check precedes
the budget loop of \eqref{eq:budget}, making it a floor the budget cannot trade
away.

\subsection{Mechanism 1 --- aggregate budget constraint}

\begin{align}
p_i(\round) &= \sigma\!\bigl(c_\lambda (\Teff_i(\round) - \tfrac{1}{2})\bigr),
\qquad \sigma(x)=\bigl(1+e^{-x}\bigr)^{-1},
\label{eq:posterior}\\[4pt]
\gamma(\round) &= \frac{\sum_{i \in \Pset} p_i(\round)}
                       {\lvert \Vset(\round)\rvert + \sum_{i \in \Pset} p_i(\round)}
\ \le\ \gamma_{\mathrm{budget}}.
\label{eq:budget}
\end{align}

While \eqref{eq:budget} is violated, the lowest-trust promoted client is moved
to $\Vset$:
\begin{equation}
i^\star = \argmin_{i \in \Pset} \Teff_i(\round),
\qquad
\Pset \leftarrow \Pset \setminus \{i^\star\},
\quad
\Vset \leftarrow \Vset \cup \{i^\star\}.
\label{eq:demote}
\end{equation}

Since demotion increases $\lvert\Vset\rvert$, the loop is monotone decreasing in
$\gamma$ and terminates. Because the ordering is ascending in trust and Tier~2
lies below $\theta_{\mathrm{high}}$ by construction, Tier~2 clients are demoted
before any Tier~3 client --- which is the sense in which Tier~2 is
\emph{budget-limited}.

Note that \eqref{eq:budget} bounds aggregation \emph{weight}, not the number of
clients and not the \emph{magnitude} of what they contribute; the latter
requires \eqref{eq:clip}.

\subsection{Trust and counter updates}

\begin{equation}
T_i(\round+1)=
\begin{cases}
\alpha\, T_i(\round) + (1-\alpha)\,\varphi_i(\round), & i \in \Vset(\round),\\[2pt]
T_i(\round), & i \in \Pset(\round).
\end{cases}
\label{eq:trust}
\end{equation}

\changed{The promoted branch was $T_i(\round+1)=\alpha T_i(\round)$. Under that
rule, a client verified a fraction $f$ of the rounds it appears in has EMA fixed
point $T^\star = f\varphi$ --- the \emph{verification rate}, not its behaviour.
Since promotion reduces $f$, trust was self-limiting: measured $T^\star \approx
0.29$ against $\theta_{\mathrm{high}}=0.7$, making Tier~3 unreachable at any
round count. Not being observed is not evidence of misbehaviour; it is a reason
to observe again, which \eqref{eq:stale} now expresses. The fixed point is now
$T^\star=\varphi$.}

\begin{align}
i \in \Vset:&\quad a_i \leftarrow 0,\qquad \ell_i \leftarrow \round,
\label{eq:reset}\\
i \in \Pset:&\quad a_i \leftarrow a_i + 1,
\label{eq:advance}\\[4pt]
c_i \leftarrow &\
\begin{cases}
c_i + 1, & i \in \Vset \ \wedge\ i \notin \Oset \ \wedge\ \varphi_i > 0.8,\\
0, & \text{otherwise.}
\end{cases}
\label{eq:streak}
\end{align}

The outlier condition in \eqref{eq:streak} is the detector's own verdict, not a
score cutoff: under \eqref{eq:phi} a client marginally past the anomaly
threshold still scores above $0.8$, so a flagged client would otherwise
continue accumulating the streak that Tier~3 requires.

% =============================================================================
\section{Layer 2 --- Ephemeral Structured Projection}

Per-round, per-block projection matrices are derived from a secret key, so no
client can anticipate the subspace its update will be examined in:
\begin{align}
\mathrm{seed}_m(\round) &= \mathrm{SHA\text{-}256}\bigl(\text{key} \,\|\, \round \,\|\, m\bigr),
\label{eq:seed}\\[2pt]
R_m(\round) &\sim \mathcal{N}\!\left(0, \tfrac{1}{k_m}\right)^{k_m \times d_m},
\qquad \textstyle\sum_m k_m = k,
\label{eq:projmat}\\[2pt]
\gproj_i^{(m)} &= R_m(\round)\, g_i^{(m)}.
\label{eq:project}
\end{align}

By the Johnson--Lindenstrauss lemma pairwise distances are preserved up to
$(1\pm\epsilon)$ with $k_m = O(\epsilon^{-2}\log N)$, so detection may run on
sketches rather than on full parameter blocks.

\medskip
\noindent
\textbf{Allocation caveat.} The implementation allocates $k_m$ proportionally to
$d_m$ with a floor of $1$. For the evaluated model at $k=150$ this yields
$k_m=1$ for six of ten blocks, where \eqref{eq:z} is a single squared Gaussian
and therefore extremely heavy-tailed. This is the source of the false-positive
behaviour discussed in \S\ref{sec:bvd}.

% =============================================================================
\section{Layer 3 --- Detection and Aggregation}
\label{sec:bvd}

\subsection{Block-variance outlier detection}

\begin{align}
\bar{g}^{(m)} &= \operatorname*{median}_{i \in \Vset} \gproj_i^{(m)}
\qquad\text{(coordinate-wise)},
\label{eq:center}\\[2pt]
d_i^{(m)} &= \bigl\lVert \gproj_i^{(m)} - \bar{g}^{(m)} \bigr\rVert_2^2,
\label{eq:dist}\\[2pt]
z_i^{(m)} &= \frac{d_i^{(m)}}{\sigsq_m^2 + \varepsilon},
\qquad
Z_i = \max_m z_i^{(m)},
\label{eq:z}\\[2pt]
\Oset(\round) &= \{\, i \in \Vset : Z_i > \tau_z \,\},
\qquad
\Lset(\round) = \Vset(\round) \setminus \Oset(\round),
\label{eq:outlier}\\[2pt]
\sigsq_m^2 &\leftarrow \alpha_\sigma\, \sigsq_m^2
+ (1-\alpha_\sigma)\,\frac{1}{\lvert\Lset\rvert}\sum_{i \in \Lset} d_i^{(m)}.
\label{eq:sigma}
\end{align}

\medskip
\noindent
\textbf{Known defects in \eqref{eq:z}--\eqref{eq:sigma}.} These are stated
because the empirical false-positive rate reaches $26\%$ on IID data with no
adversary present, and the cause is structural rather than a matter of tuning
$\tau_z$:
\begin{enumerate}\itemsep2pt
\item $Z_i$ is a ratio of a squared distance to a \emph{mean} squared distance,
      not a standardised score --- no dispersion estimate appears in the
      denominator, so $\tau_z$ admits no interpretation as a tail probability.
\item $\sigsq_m^2$ in \eqref{eq:sigma} is estimated only over $\Lset$, i.e. over
      clients already accepted. This truncates the upper tail of the quantity
      being estimated and biases it downward on every round.
\item $\max_m$ in \eqref{eq:z} carries no multiplicity correction, so the
      detector becomes stricter as blocks are added, and the block with the
      smallest $\sigsq_m^2$ determines every verdict. Measured divergence
      between the smallest and largest $\sigsq_m^2$ reaches $4\times10^3$ over
      $100$ rounds.
\end{enumerate}

\subsection{Behaviour score}

\begin{equation}
\varphi_i(\round)=
\begin{cases}
1, & Z_i \le \tau_z,\\[2pt]
\max\!\left(0,\ 1 - \dfrac{Z_i - \tau_z}{\tau_z}\right), & Z_i > \tau_z.
\end{cases}
\label{eq:phi}
\end{equation}

\changed{Previously $\varphi_i = 1 - d_i / \max_j d_j$, i.e.\ normalised by the
round's \emph{worst} client. One client therefore scored $0$ in every round,
including all-honest cohorts, and the cohort mean was pinned near
$1 - \bar{d}/\max_j d_j \approx 0.3$. Since \eqref{eq:trust} converges to
$\varphi$, this was a second independent mechanism capping trust below
$\theta_{\mathrm{high}}$. The score is now absolute, reusing the detector's own
anomaly threshold rather than introducing a further constant.}

\subsection{Magnitude bound on unverified updates}

\begin{align}
c &= \operatorname*{median}_{j \in \Vset} g_j
\qquad\text{(coordinate-wise)},
\label{eq:consensus}\\[2pt]
\tau &= \kappa \cdot \operatorname*{median}_{j \in \Vset} \lVert g_j - c \rVert_2,
\label{eq:radius}\\[2pt]
g_i &\leftarrow c + (g_i - c)\cdot\min\!\left(1,\ \frac{\tau}{\lVert g_i - c\rVert_2}\right),
\qquad i \in \Pset.
\label{eq:clip}
\end{align}

The radius is self-calibrating: it tracks the verified cohort's own spread, so
no absolute scale constant is required, and $\kappa=1$ states exactly that a
promoted client may deviate as much as a typical verified client does. Only
magnitude is affected --- direction is preserved and clients inside the ball are
returned unchanged. Verified clients are exempt, having already passed
\eqref{eq:outlier}.

\medskip
\noindent
\textbf{Why \eqref{eq:budget} alone is insufficient.} The budget bounds the
weight unverified clients hold, not the norm of what they submit. A measured
configuration with $\gamma = 0.114$ produced an aggregate norm of order $10^6$;
the two bounds constrain orthogonal quantities and neither subsumes the other.

\subsection{Aggregation}

\begin{align}
v_i &= \max(0.05,\ \varphi_i)\cdot n_i, \qquad i \in \Lset,
\label{eq:vweight}\\[2pt]
w_i &= p_i(\round)\cdot n_i, \qquad\qquad\ \ i \in \Pset,
\label{eq:pweight}\\[2pt]
Z(\round) &= \sum_{i \in \Lset} v_i + \sum_{i \in \Pset} w_i,
\label{eq:norm}\\[4pt]
g(\round+1) &= \frac{1}{Z(\round)}
\left( \sum_{i \in \Lset} v_i\, g_i + \sum_{i \in \Pset} w_i\, g_i \right).
\label{eq:aggregate}
\end{align}

\changed{The factor $n_i$ in \eqref{eq:vweight}--\eqref{eq:pweight} was absent:
weighting derived from trust alone. Under a Dirichlet split at $\alpha=0.3$ the
measured spread of client dataset sizes is $6.5\times$ (811 to 5261 samples), so
a client holding $811$ samples carried the same weight as one holding $5261$.
The trust term is retained as a multiplier, so unverified clients are still
discounted --- but relative to a correct FedAvg base weight rather than in place
of it. Aggregation is invariant to a common rescaling of all weights, which is
what makes this safe: clients submit full parameters, so a change in overall
scale would corrupt the global model rather than re-balance it.}

% =============================================================================
\section{Summary of divergences from the manuscript}

\begin{center}
\begin{tabular}{@{}llp{6.6cm}@{}}
\toprule
\textbf{Eq.} & \textbf{Quantity} & \textbf{Change} \\
\midrule
\eqref{eq:trust} & Trust update &
Promotion no longer decays trust; staleness counters
\eqref{eq:stale} carry that role. Fixed point moves from $f\varphi$ to
$\varphi$, making Tier~3 reachable. \\[4pt]
\eqref{eq:phi} & Behaviour score &
Absolute rather than relative to the round's worst client. \\[4pt]
\eqref{eq:vweight}--\eqref{eq:pweight} & Aggregation weights &
Dataset size $n_i$ restored (standard FedAvg). \\[4pt]
\eqref{eq:stale}, \eqref{eq:assign} & Staleness expiry &
New. Bounds the interval a promoted client may go unverified, in both
appearances and wall-clock rounds. \\
\bottomrule
\end{tabular}
\end{center}

\medskip
\noindent
Equations \eqref{eq:z}--\eqref{eq:sigma} are reported as implemented and are
known to be defective; they are the subject of ongoing revision and should not
be presented as a validated component.

\end{document}
 (then delete the preamble
%  and \begin{document}/\end{document}).
% ============================================================================
\documentclass[11pt]{article}

\usepackage[margin=1in]{geometry}
\usepackage{amsmath,amssymb,amsthm}
\DeclareMathOperator*{\argmin}{arg\,min}  % not provided by amsmath
\usepackage{booktabs}
\usepackage{xcolor}
\usepackage[colorlinks=true,linkcolor=black,citecolor=black]{hyperref}

% Divergences from the original formulation are flagged, not hidden.
\newcommand{\changed}[1]{%
  \par\vspace{2pt}\noindent
  \textcolor{red!65!black}{\small\textbf{Differs from manuscript:} #1}\par\vspace{2pt}}

% ---- notation shorthands -----------------------------------------------------
\newcommand{\round}{r}
\newcommand{\Vset}{\mathcal{V}}          % verified
\newcommand{\Pset}{\mathcal{P}}          % promoted
\newcommand{\Dset}{\mathcal{D}}          % decoys
\newcommand{\Lset}{\mathcal{L}}          % inliers
\newcommand{\Oset}{\mathcal{O}}          % outliers
\newcommand{\Teff}{\widetilde{T}}        % effective (ramp-capped) trust
\newcommand{\gproj}{\hat{g}}             % projected update
\newcommand{\sigsq}{\hat{\sigma}}
\newcommand{\1}[1]{\mathbf{1}\!\left[#1\right]}  % amssymb \mathbb covers uppercase only

\title{\textbf{TAVS--ESP: Formal Specification}\\[2pt]
       \large Trust-Adaptive Verification Scheduling with Ephemeral Structured Projection}
\author{}
\date{}

\begin{document}
\maketitle
\vspace{-3em}

% =============================================================================
\section{Notation}

\begin{center}
\begin{tabular}{@{}ll@{\qquad}ll@{}}
\toprule
$\round$ & round index & $i$ & client index \\
$m$ & parameter block index & $g_i$ & parameters submitted by client $i$ \\
$n_i$ & size of client $i$'s dataset & $d_m$ & dimension of block $m$ \\
$T_i(\round)$ & trust score & $\Teff_i(\round)$ & effective (ramp-capped) trust \\
$c_i$ & consecutive clean verifications & $a_i$ & appearances since last verified \\
$\ell_i$ & round of last verification & $r_0^i$ & round client $i$ joined \\
$\Vset(\round)$ & verified set & $\Pset(\round)$ & promoted set \\
$\Dset(\round)\subseteq\Vset(\round)$ & decoys & $\Lset(\round)$ & inliers \\
\bottomrule
\end{tabular}
\end{center}

\medskip
\noindent
Parameters: $\theta_{\mathrm{low}},\theta_{\mathrm{high}}$ (tier thresholds),
$\alpha$ (trust EMA), $\tau_{\mathrm{ramp}}$, $k_{\mathrm{trust}}$,
$p_{\mathrm{decoy}}$, $c_\lambda$, $\gamma_{\mathrm{budget}}$,
$s^{\mathrm{app}}_{\max}$, $s^{\mathrm{rnd}}_{\max}$,
$\tau_z$, $\alpha_\sigma$, $\kappa$ (clip factor), $k$ (projection budget).

% =============================================================================
\section{Layer 1 --- Trust-Adaptive Verification Scheduling}

\subsection{Effective trust}

A newly joined client cannot accrue influence quickly, regardless of behaviour
(Sybil resistance):
\begin{align}
T^{\max}_i(\round) &= 1 - \exp\!\left(-\frac{\round - r_0^i}{\tau_{\mathrm{ramp}}}\right),
\label{eq:ramp}\\[2pt]
\Teff_i(\round)    &= \min\!\bigl(T_i(\round),\, T^{\max}_i(\round)\bigr).
\label{eq:efftrust}
\end{align}

\subsection{Tier assignment}

\begin{equation}
\mathrm{tier}_i(\round)=
\begin{cases}
1, & \Teff_i(\round) < \theta_{\mathrm{low}},\\[2pt]
3, & \Teff_i(\round) \ge \theta_{\mathrm{high}} \ \wedge\  c_i \ge k_{\mathrm{trust}},\\[2pt]
2, & \text{otherwise.}
\end{cases}
\label{eq:tier}
\end{equation}

Tier~3 requires \emph{both} conditions: a trust threshold alone would let a
single favourable round grant promotion.

\subsection{Staleness}

Trust states how well a client behaved when last observed; staleness states
whether that observation is still current. The two are tracked separately.
\begin{equation}
\mathrm{stale}_i(\round)=
\1{a_i \ge s^{\mathrm{app}}_{\max}}
\ \vee\
\1{\round - \ell_i \ge s^{\mathrm{rnd}}_{\max}}
\ \vee\
\1{\text{$i$ never verified}}.
\label{eq:stale}
\end{equation}

Both caps are required: they measure different quantities. At $40\%$
participation, $s^{\mathrm{app}}_{\max}=4$ appearances corresponds to roughly
ten wall-clock rounds, so neither bound implies the other. A client that has
never been verified is stale by definition and cannot be promoted on no
evidence.

\subsection{Assignment rule}

\begin{equation}
i \longmapsto
\begin{cases}
\Vset, & \mathrm{stale}_i(\round) \quad \text{(hard expiry, outranks tier)},\\[2pt]
\Vset, & \mathrm{tier}_i = 1,\\[2pt]
\Vset \ \text{and} \ \Dset, & \mathrm{tier}_i = 3 \ \wedge\ \mathrm{CSPRNG}(\mathrm{key},i,\round) < p_{\mathrm{decoy}},\\[2pt]
\Pset, & \mathrm{tier}_i = 3 \ \wedge\ \text{otherwise},\\[2pt]
\Pset, & \mathrm{tier}_i = 2 \quad \text{(tentative; see \eqref{eq:budget})}.
\end{cases}
\label{eq:assign}
\end{equation}

Decoys are verified server-side while the client is told it was promoted, so no
client can determine whether it is being observed. The expiry check precedes
the budget loop of \eqref{eq:budget}, making it a floor the budget cannot trade
away.

\subsection{Mechanism 1 --- aggregate budget constraint}

\begin{align}
p_i(\round) &= \sigma\!\bigl(c_\lambda (\Teff_i(\round) - \tfrac{1}{2})\bigr),
\qquad \sigma(x)=\bigl(1+e^{-x}\bigr)^{-1},
\label{eq:posterior}\\[4pt]
\gamma(\round) &= \frac{\sum_{i \in \Pset} p_i(\round)}
                       {\lvert \Vset(\round)\rvert + \sum_{i \in \Pset} p_i(\round)}
\ \le\ \gamma_{\mathrm{budget}}.
\label{eq:budget}
\end{align}

While \eqref{eq:budget} is violated, the lowest-trust promoted client is moved
to $\Vset$:
\begin{equation}
i^\star = \argmin_{i \in \Pset} \Teff_i(\round),
\qquad
\Pset \leftarrow \Pset \setminus \{i^\star\},
\quad
\Vset \leftarrow \Vset \cup \{i^\star\}.
\label{eq:demote}
\end{equation}

Since demotion increases $\lvert\Vset\rvert$, the loop is monotone decreasing in
$\gamma$ and terminates. Because the ordering is ascending in trust and Tier~2
lies below $\theta_{\mathrm{high}}$ by construction, Tier~2 clients are demoted
before any Tier~3 client --- which is the sense in which Tier~2 is
\emph{budget-limited}.

Note that \eqref{eq:budget} bounds aggregation \emph{weight}, not the number of
clients and not the \emph{magnitude} of what they contribute; the latter
requires \eqref{eq:clip}.

\subsection{Trust and counter updates}

\begin{equation}
T_i(\round+1)=
\begin{cases}
\alpha\, T_i(\round) + (1-\alpha)\,\varphi_i(\round), & i \in \Vset(\round),\\[2pt]
T_i(\round), & i \in \Pset(\round).
\end{cases}
\label{eq:trust}
\end{equation}

\changed{The promoted branch was $T_i(\round+1)=\alpha T_i(\round)$. Under that
rule, a client verified a fraction $f$ of the rounds it appears in has EMA fixed
point $T^\star = f\varphi$ --- the \emph{verification rate}, not its behaviour.
Since promotion reduces $f$, trust was self-limiting: measured $T^\star \approx
0.29$ against $\theta_{\mathrm{high}}=0.7$, making Tier~3 unreachable at any
round count. Not being observed is not evidence of misbehaviour; it is a reason
to observe again, which \eqref{eq:stale} now expresses. The fixed point is now
$T^\star=\varphi$.}

\begin{align}
i \in \Vset:&\quad a_i \leftarrow 0,\qquad \ell_i \leftarrow \round,
\label{eq:reset}\\
i \in \Pset:&\quad a_i \leftarrow a_i + 1,
\label{eq:advance}\\[4pt]
c_i \leftarrow &\
\begin{cases}
c_i + 1, & i \in \Vset \ \wedge\ i \notin \Oset \ \wedge\ \varphi_i > 0.8,\\
0, & \text{otherwise.}
\end{cases}
\label{eq:streak}
\end{align}

The outlier condition in \eqref{eq:streak} is the detector's own verdict, not a
score cutoff: under \eqref{eq:phi} a client marginally past the anomaly
threshold still scores above $0.8$, so a flagged client would otherwise
continue accumulating the streak that Tier~3 requires.

% =============================================================================
\section{Layer 2 --- Ephemeral Structured Projection}

Per-round, per-block projection matrices are derived from a secret key, so no
client can anticipate the subspace its update will be examined in:
\begin{align}
\mathrm{seed}_m(\round) &= \mathrm{SHA\text{-}256}\bigl(\text{key} \,\|\, \round \,\|\, m\bigr),
\label{eq:seed}\\[2pt]
R_m(\round) &\sim \mathcal{N}\!\left(0, \tfrac{1}{k_m}\right)^{k_m \times d_m},
\qquad \textstyle\sum_m k_m = k,
\label{eq:projmat}\\[2pt]
\gproj_i^{(m)} &= R_m(\round)\, g_i^{(m)}.
\label{eq:project}
\end{align}

By the Johnson--Lindenstrauss lemma pairwise distances are preserved up to
$(1\pm\epsilon)$ with $k_m = O(\epsilon^{-2}\log N)$, so detection may run on
sketches rather than on full parameter blocks.

\medskip
\noindent
\textbf{Allocation caveat.} The implementation allocates $k_m$ proportionally to
$d_m$ with a floor of $1$. For the evaluated model at $k=150$ this yields
$k_m=1$ for six of ten blocks, where \eqref{eq:z} is a single squared Gaussian
and therefore extremely heavy-tailed. This is the source of the false-positive
behaviour discussed in \S\ref{sec:bvd}.

% =============================================================================
\section{Layer 3 --- Detection and Aggregation}
\label{sec:bvd}

\subsection{Block-variance outlier detection}

\begin{align}
\bar{g}^{(m)} &= \operatorname*{median}_{i \in \Vset} \gproj_i^{(m)}
\qquad\text{(coordinate-wise)},
\label{eq:center}\\[2pt]
d_i^{(m)} &= \bigl\lVert \gproj_i^{(m)} - \bar{g}^{(m)} \bigr\rVert_2^2,
\label{eq:dist}\\[2pt]
z_i^{(m)} &= \frac{d_i^{(m)}}{\sigsq_m^2 + \varepsilon},
\qquad
Z_i = \max_m z_i^{(m)},
\label{eq:z}\\[2pt]
\Oset(\round) &= \{\, i \in \Vset : Z_i > \tau_z \,\},
\qquad
\Lset(\round) = \Vset(\round) \setminus \Oset(\round),
\label{eq:outlier}\\[2pt]
\sigsq_m^2 &\leftarrow \alpha_\sigma\, \sigsq_m^2
+ (1-\alpha_\sigma)\,\frac{1}{\lvert\Lset\rvert}\sum_{i \in \Lset} d_i^{(m)}.
\label{eq:sigma}
\end{align}

\medskip
\noindent
\textbf{Known defects in \eqref{eq:z}--\eqref{eq:sigma}.} These are stated
because the empirical false-positive rate reaches $26\%$ on IID data with no
adversary present, and the cause is structural rather than a matter of tuning
$\tau_z$:
\begin{enumerate}\itemsep2pt
\item $Z_i$ is a ratio of a squared distance to a \emph{mean} squared distance,
      not a standardised score --- no dispersion estimate appears in the
      denominator, so $\tau_z$ admits no interpretation as a tail probability.
\item $\sigsq_m^2$ in \eqref{eq:sigma} is estimated only over $\Lset$, i.e. over
      clients already accepted. This truncates the upper tail of the quantity
      being estimated and biases it downward on every round.
\item $\max_m$ in \eqref{eq:z} carries no multiplicity correction, so the
      detector becomes stricter as blocks are added, and the block with the
      smallest $\sigsq_m^2$ determines every verdict. Measured divergence
      between the smallest and largest $\sigsq_m^2$ reaches $4\times10^3$ over
      $100$ rounds.
\end{enumerate}

\subsection{Behaviour score}

\begin{equation}
\varphi_i(\round)=
\begin{cases}
1, & Z_i \le \tau_z,\\[2pt]
\max\!\left(0,\ 1 - \dfrac{Z_i - \tau_z}{\tau_z}\right), & Z_i > \tau_z.
\end{cases}
\label{eq:phi}
\end{equation}

\changed{Previously $\varphi_i = 1 - d_i / \max_j d_j$, i.e.\ normalised by the
round's \emph{worst} client. One client therefore scored $0$ in every round,
including all-honest cohorts, and the cohort mean was pinned near
$1 - \bar{d}/\max_j d_j \approx 0.3$. Since \eqref{eq:trust} converges to
$\varphi$, this was a second independent mechanism capping trust below
$\theta_{\mathrm{high}}$. The score is now absolute, reusing the detector's own
anomaly threshold rather than introducing a further constant.}

\subsection{Magnitude bound on unverified updates}

\begin{align}
c &= \operatorname*{median}_{j \in \Vset} g_j
\qquad\text{(coordinate-wise)},
\label{eq:consensus}\\[2pt]
\tau &= \kappa \cdot \operatorname*{median}_{j \in \Vset} \lVert g_j - c \rVert_2,
\label{eq:radius}\\[2pt]
g_i &\leftarrow c + (g_i - c)\cdot\min\!\left(1,\ \frac{\tau}{\lVert g_i - c\rVert_2}\right),
\qquad i \in \Pset.
\label{eq:clip}
\end{align}

The radius is self-calibrating: it tracks the verified cohort's own spread, so
no absolute scale constant is required, and $\kappa=1$ states exactly that a
promoted client may deviate as much as a typical verified client does. Only
magnitude is affected --- direction is preserved and clients inside the ball are
returned unchanged. Verified clients are exempt, having already passed
\eqref{eq:outlier}.

\medskip
\noindent
\textbf{Why \eqref{eq:budget} alone is insufficient.} The budget bounds the
weight unverified clients hold, not the norm of what they submit. A measured
configuration with $\gamma = 0.114$ produced an aggregate norm of order $10^6$;
the two bounds constrain orthogonal quantities and neither subsumes the other.

\subsection{Aggregation}

\begin{align}
v_i &= \max(0.05,\ \varphi_i)\cdot n_i, \qquad i \in \Lset,
\label{eq:vweight}\\[2pt]
w_i &= p_i(\round)\cdot n_i, \qquad\qquad\ \ i \in \Pset,
\label{eq:pweight}\\[2pt]
Z(\round) &= \sum_{i \in \Lset} v_i + \sum_{i \in \Pset} w_i,
\label{eq:norm}\\[4pt]
g(\round+1) &= \frac{1}{Z(\round)}
\left( \sum_{i \in \Lset} v_i\, g_i + \sum_{i \in \Pset} w_i\, g_i \right).
\label{eq:aggregate}
\end{align}

\changed{The factor $n_i$ in \eqref{eq:vweight}--\eqref{eq:pweight} was absent:
weighting derived from trust alone. Under a Dirichlet split at $\alpha=0.3$ the
measured spread of client dataset sizes is $6.5\times$ (811 to 5261 samples), so
a client holding $811$ samples carried the same weight as one holding $5261$.
The trust term is retained as a multiplier, so unverified clients are still
discounted --- but relative to a correct FedAvg base weight rather than in place
of it. Aggregation is invariant to a common rescaling of all weights, which is
what makes this safe: clients submit full parameters, so a change in overall
scale would corrupt the global model rather than re-balance it.}

% =============================================================================
\section{Summary of divergences from the manuscript}

\begin{center}
\begin{tabular}{@{}llp{6.6cm}@{}}
\toprule
\textbf{Eq.} & \textbf{Quantity} & \textbf{Change} \\
\midrule
\eqref{eq:trust} & Trust update &
Promotion no longer decays trust; staleness counters
\eqref{eq:stale} carry that role. Fixed point moves from $f\varphi$ to
$\varphi$, making Tier~3 reachable. \\[4pt]
\eqref{eq:phi} & Behaviour score &
Absolute rather than relative to the round's worst client. \\[4pt]
\eqref{eq:vweight}--\eqref{eq:pweight} & Aggregation weights &
Dataset size $n_i$ restored (standard FedAvg). \\[4pt]
\eqref{eq:stale}, \eqref{eq:assign} & Staleness expiry &
New. Bounds the interval a promoted client may go unverified, in both
appearances and wall-clock rounds. \\
\bottomrule
\end{tabular}
\end{center}

\medskip
\noindent
Equations \eqref{eq:z}--\eqref{eq:sigma} are reported as implemented and are
known to be defective; they are the subject of ongoing revision and should not
be presented as a validated component.

\end{document}
 (then delete the preamble
%  and \begin{document}/\end{document}).
% ============================================================================
\documentclass[11pt]{article}

\usepackage[margin=1in]{geometry}
\usepackage{amsmath,amssymb,amsthm}
\DeclareMathOperator*{\argmin}{arg\,min}  % not provided by amsmath
\usepackage{booktabs}
\usepackage{xcolor}
\usepackage[colorlinks=true,linkcolor=black,citecolor=black]{hyperref}

% Divergences from the original formulation are flagged, not hidden.
\newcommand{\changed}[1]{%
  \par\vspace{2pt}\noindent
  \textcolor{red!65!black}{\small\textbf{Differs from manuscript:} #1}\par\vspace{2pt}}

% ---- notation shorthands -----------------------------------------------------
\newcommand{\round}{r}
\newcommand{\Vset}{\mathcal{V}}          % verified
\newcommand{\Pset}{\mathcal{P}}          % promoted
\newcommand{\Dset}{\mathcal{D}}          % decoys
\newcommand{\Lset}{\mathcal{L}}          % inliers
\newcommand{\Oset}{\mathcal{O}}          % outliers
\newcommand{\Teff}{\widetilde{T}}        % effective (ramp-capped) trust
\newcommand{\gproj}{\hat{g}}             % projected update
\newcommand{\sigsq}{\hat{\sigma}}
\newcommand{\1}[1]{\mathbf{1}\!\left[#1\right]}  % amssymb \mathbb covers uppercase only

\title{\textbf{TAVS--ESP: Formal Specification}\\[2pt]
       \large Trust-Adaptive Verification Scheduling with Ephemeral Structured Projection}
\author{}
\date{}

\begin{document}
\maketitle
\vspace{-3em}

% =============================================================================
\section{Notation}

\begin{center}
\begin{tabular}{@{}ll@{\qquad}ll@{}}
\toprule
$\round$ & round index & $i$ & client index \\
$m$ & parameter block index & $g_i$ & parameters submitted by client $i$ \\
$n_i$ & size of client $i$'s dataset & $d_m$ & dimension of block $m$ \\
$T_i(\round)$ & trust score & $\Teff_i(\round)$ & effective (ramp-capped) trust \\
$c_i$ & consecutive clean verifications & $a_i$ & appearances since last verified \\
$\ell_i$ & round of last verification & $r_0^i$ & round client $i$ joined \\
$\Vset(\round)$ & verified set & $\Pset(\round)$ & promoted set \\
$\Dset(\round)\subseteq\Vset(\round)$ & decoys & $\Lset(\round)$ & inliers \\
\bottomrule
\end{tabular}
\end{center}

\medskip
\noindent
Parameters: $\theta_{\mathrm{low}},\theta_{\mathrm{high}}$ (tier thresholds),
$\alpha$ (trust EMA), $\tau_{\mathrm{ramp}}$, $k_{\mathrm{trust}}$,
$p_{\mathrm{decoy}}$, $c_\lambda$, $\gamma_{\mathrm{budget}}$,
$s^{\mathrm{app}}_{\max}$, $s^{\mathrm{rnd}}_{\max}$,
$\tau_z$, $\alpha_\sigma$, $\kappa$ (clip factor), $k$ (projection budget).

% =============================================================================
\section{Layer 1 --- Trust-Adaptive Verification Scheduling}

\subsection{Effective trust}

A newly joined client cannot accrue influence quickly, regardless of behaviour
(Sybil resistance):
\begin{align}
T^{\max}_i(\round) &= 1 - \exp\!\left(-\frac{\round - r_0^i}{\tau_{\mathrm{ramp}}}\right),
\label{eq:ramp}\\[2pt]
\Teff_i(\round)    &= \min\!\bigl(T_i(\round),\, T^{\max}_i(\round)\bigr).
\label{eq:efftrust}
\end{align}

\subsection{Tier assignment}

\begin{equation}
\mathrm{tier}_i(\round)=
\begin{cases}
1, & \Teff_i(\round) < \theta_{\mathrm{low}},\\[2pt]
3, & \Teff_i(\round) \ge \theta_{\mathrm{high}} \ \wedge\  c_i \ge k_{\mathrm{trust}},\\[2pt]
2, & \text{otherwise.}
\end{cases}
\label{eq:tier}
\end{equation}

Tier~3 requires \emph{both} conditions: a trust threshold alone would let a
single favourable round grant promotion.

\subsection{Staleness}

Trust states how well a client behaved when last observed; staleness states
whether that observation is still current. The two are tracked separately.
\begin{equation}
\mathrm{stale}_i(\round)=
\1{a_i \ge s^{\mathrm{app}}_{\max}}
\ \vee\
\1{\round - \ell_i \ge s^{\mathrm{rnd}}_{\max}}
\ \vee\
\1{\text{$i$ never verified}}.
\label{eq:stale}
\end{equation}

Both caps are required: they measure different quantities. At $40\%$
participation, $s^{\mathrm{app}}_{\max}=4$ appearances corresponds to roughly
ten wall-clock rounds, so neither bound implies the other. A client that has
never been verified is stale by definition and cannot be promoted on no
evidence.

\subsection{Assignment rule}

\begin{equation}
i \longmapsto
\begin{cases}
\Vset, & \mathrm{stale}_i(\round) \quad \text{(hard expiry, outranks tier)},\\[2pt]
\Vset, & \mathrm{tier}_i = 1,\\[2pt]
\Vset \ \text{and} \ \Dset, & \mathrm{tier}_i = 3 \ \wedge\ \mathrm{CSPRNG}(\mathrm{key},i,\round) < p_{\mathrm{decoy}},\\[2pt]
\Pset, & \mathrm{tier}_i = 3 \ \wedge\ \text{otherwise},\\[2pt]
\Pset, & \mathrm{tier}_i = 2 \quad \text{(tentative; see \eqref{eq:budget})}.
\end{cases}
\label{eq:assign}
\end{equation}

Decoys are verified server-side while the client is told it was promoted, so no
client can determine whether it is being observed. The expiry check precedes
the budget loop of \eqref{eq:budget}, making it a floor the budget cannot trade
away.

\subsection{Mechanism 1 --- aggregate budget constraint}

\begin{align}
p_i(\round) &= \sigma\!\bigl(c_\lambda (\Teff_i(\round) - \tfrac{1}{2})\bigr),
\qquad \sigma(x)=\bigl(1+e^{-x}\bigr)^{-1},
\label{eq:posterior}\\[4pt]
\gamma(\round) &= \frac{\sum_{i \in \Pset} p_i(\round)}
                       {\lvert \Vset(\round)\rvert + \sum_{i \in \Pset} p_i(\round)}
\ \le\ \gamma_{\mathrm{budget}}.
\label{eq:budget}
\end{align}

While \eqref{eq:budget} is violated, the lowest-trust promoted client is moved
to $\Vset$:
\begin{equation}
i^\star = \argmin_{i \in \Pset} \Teff_i(\round),
\qquad
\Pset \leftarrow \Pset \setminus \{i^\star\},
\quad
\Vset \leftarrow \Vset \cup \{i^\star\}.
\label{eq:demote}
\end{equation}

Since demotion increases $\lvert\Vset\rvert$, the loop is monotone decreasing in
$\gamma$ and terminates. Because the ordering is ascending in trust and Tier~2
lies below $\theta_{\mathrm{high}}$ by construction, Tier~2 clients are demoted
before any Tier~3 client --- which is the sense in which Tier~2 is
\emph{budget-limited}.

Note that \eqref{eq:budget} bounds aggregation \emph{weight}, not the number of
clients and not the \emph{magnitude} of what they contribute; the latter
requires \eqref{eq:clip}.

\subsection{Trust and counter updates}

\begin{equation}
T_i(\round+1)=
\begin{cases}
\alpha\, T_i(\round) + (1-\alpha)\,\varphi_i(\round), & i \in \Vset(\round),\\[2pt]
T_i(\round), & i \in \Pset(\round).
\end{cases}
\label{eq:trust}
\end{equation}

\changed{The promoted branch was $T_i(\round+1)=\alpha T_i(\round)$. Under that
rule, a client verified a fraction $f$ of the rounds it appears in has EMA fixed
point $T^\star = f\varphi$ --- the \emph{verification rate}, not its behaviour.
Since promotion reduces $f$, trust was self-limiting: measured $T^\star \approx
0.29$ against $\theta_{\mathrm{high}}=0.7$, making Tier~3 unreachable at any
round count. Not being observed is not evidence of misbehaviour; it is a reason
to observe again, which \eqref{eq:stale} now expresses. The fixed point is now
$T^\star=\varphi$.}

\begin{align}
i \in \Vset:&\quad a_i \leftarrow 0,\qquad \ell_i \leftarrow \round,
\label{eq:reset}\\
i \in \Pset:&\quad a_i \leftarrow a_i + 1,
\label{eq:advance}\\[4pt]
c_i \leftarrow &\
\begin{cases}
c_i + 1, & i \in \Vset \ \wedge\ i \notin \Oset \ \wedge\ \varphi_i > 0.8,\\
0, & \text{otherwise.}
\end{cases}
\label{eq:streak}
\end{align}

The outlier condition in \eqref{eq:streak} is the detector's own verdict, not a
score cutoff: under \eqref{eq:phi} a client marginally past the anomaly
threshold still scores above $0.8$, so a flagged client would otherwise
continue accumulating the streak that Tier~3 requires.

% =============================================================================
\section{Layer 2 --- Ephemeral Structured Projection}

Per-round, per-block projection matrices are derived from a secret key, so no
client can anticipate the subspace its update will be examined in:
\begin{align}
\mathrm{seed}_m(\round) &= \mathrm{SHA\text{-}256}\bigl(\text{key} \,\|\, \round \,\|\, m\bigr),
\label{eq:seed}\\[2pt]
R_m(\round) &\sim \mathcal{N}\!\left(0, \tfrac{1}{k_m}\right)^{k_m \times d_m},
\qquad \textstyle\sum_m k_m = k,
\label{eq:projmat}\\[2pt]
\gproj_i^{(m)} &= R_m(\round)\, g_i^{(m)}.
\label{eq:project}
\end{align}

By the Johnson--Lindenstrauss lemma pairwise distances are preserved up to
$(1\pm\epsilon)$ with $k_m = O(\epsilon^{-2}\log N)$, so detection may run on
sketches rather than on full parameter blocks.

\medskip
\noindent
\textbf{Allocation caveat.} The implementation allocates $k_m$ proportionally to
$d_m$ with a floor of $1$. For the evaluated model at $k=150$ this yields
$k_m=1$ for six of ten blocks, where \eqref{eq:z} is a single squared Gaussian
and therefore extremely heavy-tailed. This is the source of the false-positive
behaviour discussed in \S\ref{sec:bvd}.

% =============================================================================
\section{Layer 3 --- Detection and Aggregation}
\label{sec:bvd}

\subsection{Block-variance outlier detection}

\begin{align}
\bar{g}^{(m)} &= \operatorname*{median}_{i \in \Vset} \gproj_i^{(m)}
\qquad\text{(coordinate-wise)},
\label{eq:center}\\[2pt]
d_i^{(m)} &= \bigl\lVert \gproj_i^{(m)} - \bar{g}^{(m)} \bigr\rVert_2^2,
\label{eq:dist}\\[2pt]
z_i^{(m)} &= \frac{d_i^{(m)}}{\sigsq_m^2 + \varepsilon},
\qquad
Z_i = \max_m z_i^{(m)},
\label{eq:z}\\[2pt]
\Oset(\round) &= \{\, i \in \Vset : Z_i > \tau_z \,\},
\qquad
\Lset(\round) = \Vset(\round) \setminus \Oset(\round),
\label{eq:outlier}\\[2pt]
\sigsq_m^2 &\leftarrow \alpha_\sigma\, \sigsq_m^2
+ (1-\alpha_\sigma)\,\frac{1}{\lvert\Lset\rvert}\sum_{i \in \Lset} d_i^{(m)}.
\label{eq:sigma}
\end{align}

\medskip
\noindent
\textbf{Known defects in \eqref{eq:z}--\eqref{eq:sigma}.} These are stated
because the empirical false-positive rate reaches $26\%$ on IID data with no
adversary present, and the cause is structural rather than a matter of tuning
$\tau_z$:
\begin{enumerate}\itemsep2pt
\item $Z_i$ is a ratio of a squared distance to a \emph{mean} squared distance,
      not a standardised score --- no dispersion estimate appears in the
      denominator, so $\tau_z$ admits no interpretation as a tail probability.
\item $\sigsq_m^2$ in \eqref{eq:sigma} is estimated only over $\Lset$, i.e. over
      clients already accepted. This truncates the upper tail of the quantity
      being estimated and biases it downward on every round.
\item $\max_m$ in \eqref{eq:z} carries no multiplicity correction, so the
      detector becomes stricter as blocks are added, and the block with the
      smallest $\sigsq_m^2$ determines every verdict. Measured divergence
      between the smallest and largest $\sigsq_m^2$ reaches $4\times10^3$ over
      $100$ rounds.
\end{enumerate}

\subsection{Behaviour score}

\begin{equation}
\varphi_i(\round)=
\begin{cases}
1, & Z_i \le \tau_z,\\[2pt]
\max\!\left(0,\ 1 - \dfrac{Z_i - \tau_z}{\tau_z}\right), & Z_i > \tau_z.
\end{cases}
\label{eq:phi}
\end{equation}

\changed{Previously $\varphi_i = 1 - d_i / \max_j d_j$, i.e.\ normalised by the
round's \emph{worst} client. One client therefore scored $0$ in every round,
including all-honest cohorts, and the cohort mean was pinned near
$1 - \bar{d}/\max_j d_j \approx 0.3$. Since \eqref{eq:trust} converges to
$\varphi$, this was a second independent mechanism capping trust below
$\theta_{\mathrm{high}}$. The score is now absolute, reusing the detector's own
anomaly threshold rather than introducing a further constant.}

\subsection{Magnitude bound on unverified updates}

\begin{align}
c &= \operatorname*{median}_{j \in \Vset} g_j
\qquad\text{(coordinate-wise)},
\label{eq:consensus}\\[2pt]
\tau &= \kappa \cdot \operatorname*{median}_{j \in \Vset} \lVert g_j - c \rVert_2,
\label{eq:radius}\\[2pt]
g_i &\leftarrow c + (g_i - c)\cdot\min\!\left(1,\ \frac{\tau}{\lVert g_i - c\rVert_2}\right),
\qquad i \in \Pset.
\label{eq:clip}
\end{align}

The radius is self-calibrating: it tracks the verified cohort's own spread, so
no absolute scale constant is required, and $\kappa=1$ states exactly that a
promoted client may deviate as much as a typical verified client does. Only
magnitude is affected --- direction is preserved and clients inside the ball are
returned unchanged. Verified clients are exempt, having already passed
\eqref{eq:outlier}.

\medskip
\noindent
\textbf{Why \eqref{eq:budget} alone is insufficient.} The budget bounds the
weight unverified clients hold, not the norm of what they submit. A measured
configuration with $\gamma = 0.114$ produced an aggregate norm of order $10^6$;
the two bounds constrain orthogonal quantities and neither subsumes the other.

\subsection{Aggregation}

\begin{align}
v_i &= \max(0.05,\ \varphi_i)\cdot n_i, \qquad i \in \Lset,
\label{eq:vweight}\\[2pt]
w_i &= p_i(\round)\cdot n_i, \qquad\qquad\ \ i \in \Pset,
\label{eq:pweight}\\[2pt]
Z(\round) &= \sum_{i \in \Lset} v_i + \sum_{i \in \Pset} w_i,
\label{eq:norm}\\[4pt]
g(\round+1) &= \frac{1}{Z(\round)}
\left( \sum_{i \in \Lset} v_i\, g_i + \sum_{i \in \Pset} w_i\, g_i \right).
\label{eq:aggregate}
\end{align}

\changed{The factor $n_i$ in \eqref{eq:vweight}--\eqref{eq:pweight} was absent:
weighting derived from trust alone. Under a Dirichlet split at $\alpha=0.3$ the
measured spread of client dataset sizes is $6.5\times$ (811 to 5261 samples), so
a client holding $811$ samples carried the same weight as one holding $5261$.
The trust term is retained as a multiplier, so unverified clients are still
discounted --- but relative to a correct FedAvg base weight rather than in place
of it. Aggregation is invariant to a common rescaling of all weights, which is
what makes this safe: clients submit full parameters, so a change in overall
scale would corrupt the global model rather than re-balance it.}

% =============================================================================
\section{Summary of divergences from the manuscript}

\begin{center}
\begin{tabular}{@{}llp{6.6cm}@{}}
\toprule
\textbf{Eq.} & \textbf{Quantity} & \textbf{Change} \\
\midrule
\eqref{eq:trust} & Trust update &
Promotion no longer decays trust; staleness counters
\eqref{eq:stale} carry that role. Fixed point moves from $f\varphi$ to
$\varphi$, making Tier~3 reachable. \\[4pt]
\eqref{eq:phi} & Behaviour score &
Absolute rather than relative to the round's worst client. \\[4pt]
\eqref{eq:vweight}--\eqref{eq:pweight} & Aggregation weights &
Dataset size $n_i$ restored (standard FedAvg). \\[4pt]
\eqref{eq:stale}, \eqref{eq:assign} & Staleness expiry &
New. Bounds the interval a promoted client may go unverified, in both
appearances and wall-clock rounds. \\
\bottomrule
\end{tabular}
\end{center}

\medskip
\noindent
Equations \eqref{eq:z}--\eqref{eq:sigma} are reported as implemented and are
known to be defective; they are the subject of ongoing revision and should not
be presented as a validated component.

\end{document}
 (then delete the preamble
%  and \begin{document}/\end{document}).
% ============================================================================
\documentclass[11pt]{article}

\usepackage[margin=1in]{geometry}
\usepackage{amsmath,amssymb,amsthm}
\DeclareMathOperator*{\argmin}{arg\,min}  % not provided by amsmath
\usepackage{booktabs}
\usepackage{xcolor}
\usepackage[colorlinks=true,linkcolor=black,citecolor=black]{hyperref}

% Divergences from the original formulation are flagged, not hidden.
\newcommand{\changed}[1]{%
  \par\vspace{2pt}\noindent
  \textcolor{red!65!black}{\small\textbf{Differs from manuscript:} #1}\par\vspace{2pt}}

% ---- notation shorthands -----------------------------------------------------
\newcommand{\round}{r}
\newcommand{\Vset}{\mathcal{V}}          % verified
\newcommand{\Pset}{\mathcal{P}}          % promoted
\newcommand{\Dset}{\mathcal{D}}          % decoys
\newcommand{\Lset}{\mathcal{L}}          % inliers
\newcommand{\Oset}{\mathcal{O}}          % outliers
\newcommand{\Teff}{\widetilde{T}}        % effective (ramp-capped) trust
\newcommand{\gproj}{\hat{g}}             % projected update
\newcommand{\sigsq}{\hat{\sigma}}
\newcommand{\1}[1]{\mathbf{1}\!\left[#1\right]}  % amssymb \mathbb covers uppercase only

\title{\textbf{TAVS--ESP: Formal Specification}\\[2pt]
       \large Trust-Adaptive Verification Scheduling with Ephemeral Structured Projection}
\author{}
\date{}

\begin{document}
\maketitle
\vspace{-3em}

% =============================================================================
\section{Notation}

\begin{center}
\begin{tabular}{@{}ll@{\qquad}ll@{}}
\toprule
$\round$ & round index & $i$ & client index \\
$m$ & parameter block index & $g_i$ & parameters submitted by client $i$ \\
$n_i$ & size of client $i$'s dataset & $d_m$ & dimension of block $m$ \\
$T_i(\round)$ & trust score & $\Teff_i(\round)$ & effective (ramp-capped) trust \\
$c_i$ & consecutive clean verifications & $a_i$ & appearances since last verified \\
$\ell_i$ & round of last verification & $r_0^i$ & round client $i$ joined \\
$\Vset(\round)$ & verified set & $\Pset(\round)$ & promoted set \\
$\Dset(\round)\subseteq\Vset(\round)$ & decoys & $\Lset(\round)$ & inliers \\
\bottomrule
\end{tabular}
\end{center}

\medskip
\noindent
Parameters: $\theta_{\mathrm{low}},\theta_{\mathrm{high}}$ (tier thresholds),
$\alpha$ (trust EMA), $\tau_{\mathrm{ramp}}$, $k_{\mathrm{trust}}$,
$p_{\mathrm{decoy}}$, $c_\lambda$, $\gamma_{\mathrm{budget}}$,
$s^{\mathrm{app}}_{\max}$, $s^{\mathrm{rnd}}_{\max}$,
$\tau_z$, $\alpha_\sigma$, $\kappa$ (clip factor), $k$ (projection budget).

% =============================================================================
\section{Layer 1 --- Trust-Adaptive Verification Scheduling}

\subsection{Effective trust}

A newly joined client cannot accrue influence quickly, regardless of behaviour
(Sybil resistance):
\begin{align}
T^{\max}_i(\round) &= 1 - \exp\!\left(-\frac{\round - r_0^i}{\tau_{\mathrm{ramp}}}\right),
\label{eq:ramp}\\[2pt]
\Teff_i(\round)    &= \min\!\bigl(T_i(\round),\, T^{\max}_i(\round)\bigr).
\label{eq:efftrust}
\end{align}

\subsection{Tier assignment}

\begin{equation}
\mathrm{tier}_i(\round)=
\begin{cases}
1, & \Teff_i(\round) < \theta_{\mathrm{low}},\\[2pt]
3, & \Teff_i(\round) \ge \theta_{\mathrm{high}} \ \wedge\  c_i \ge k_{\mathrm{trust}},\\[2pt]
2, & \text{otherwise.}
\end{cases}
\label{eq:tier}
\end{equation}

Tier~3 requires \emph{both} conditions: a trust threshold alone would let a
single favourable round grant promotion.

\subsection{Staleness}

Trust states how well a client behaved when last observed; staleness states
whether that observation is still current. The two are tracked separately.
\begin{equation}
\mathrm{stale}_i(\round)=
\1{a_i \ge s^{\mathrm{app}}_{\max}}
\ \vee\
\1{\round - \ell_i \ge s^{\mathrm{rnd}}_{\max}}
\ \vee\
\1{\text{$i$ never verified}}.
\label{eq:stale}
\end{equation}

Both caps are required: they measure different quantities. At $40\%$
participation, $s^{\mathrm{app}}_{\max}=4$ appearances corresponds to roughly
ten wall-clock rounds, so neither bound implies the other. A client that has
never been verified is stale by definition and cannot be promoted on no
evidence.

\subsection{Assignment rule}

\begin{equation}
i \longmapsto
\begin{cases}
\Vset, & \mathrm{stale}_i(\round) \quad \text{(hard expiry, outranks tier)},\\[2pt]
\Vset, & \mathrm{tier}_i = 1,\\[2pt]
\Vset \ \text{and} \ \Dset, & \mathrm{tier}_i = 3 \ \wedge\ \mathrm{CSPRNG}(\mathrm{key},i,\round) < p_{\mathrm{decoy}},\\[2pt]
\Pset, & \mathrm{tier}_i = 3 \ \wedge\ \text{otherwise},\\[2pt]
\Pset, & \mathrm{tier}_i = 2 \quad \text{(tentative; see \eqref{eq:budget})}.
\end{cases}
\label{eq:assign}
\end{equation}

Decoys are verified server-side while the client is told it was promoted, so no
client can determine whether it is being observed. The expiry check precedes
the budget loop of \eqref{eq:budget}, making it a floor the budget cannot trade
away.

\subsection{Mechanism 1 --- aggregate budget constraint}

\begin{align}
p_i(\round) &= \sigma\!\bigl(c_\lambda (\Teff_i(\round) - \tfrac{1}{2})\bigr),
\qquad \sigma(x)=\bigl(1+e^{-x}\bigr)^{-1},
\label{eq:posterior}\\[4pt]
\gamma(\round) &= \frac{\sum_{i \in \Pset} p_i(\round)}
                       {\lvert \Vset(\round)\rvert + \sum_{i \in \Pset} p_i(\round)}
\ \le\ \gamma_{\mathrm{budget}}.
\label{eq:budget}
\end{align}

While \eqref{eq:budget} is violated, the lowest-trust promoted client is moved
to $\Vset$:
\begin{equation}
i^\star = \argmin_{i \in \Pset} \Teff_i(\round),
\qquad
\Pset \leftarrow \Pset \setminus \{i^\star\},
\quad
\Vset \leftarrow \Vset \cup \{i^\star\}.
\label{eq:demote}
\end{equation}

Since demotion increases $\lvert\Vset\rvert$, the loop is monotone decreasing in
$\gamma$ and terminates. Because the ordering is ascending in trust and Tier~2
lies below $\theta_{\mathrm{high}}$ by construction, Tier~2 clients are demoted
before any Tier~3 client --- which is the sense in which Tier~2 is
\emph{budget-limited}.

Note that \eqref{eq:budget} bounds aggregation \emph{weight}, not the number of
clients and not the \emph{magnitude} of what they contribute; the latter
requires \eqref{eq:clip}.

\subsection{Trust and counter updates}

\begin{equation}
T_i(\round+1)=
\begin{cases}
\alpha\, T_i(\round) + (1-\alpha)\,\varphi_i(\round), & i \in \Vset(\round),\\[2pt]
T_i(\round), & i \in \Pset(\round).
\end{cases}
\label{eq:trust}
\end{equation}

\changed{The promoted branch was $T_i(\round+1)=\alpha T_i(\round)$. Under that
rule, a client verified a fraction $f$ of the rounds it appears in has EMA fixed
point $T^\star = f\varphi$ --- the \emph{verification rate}, not its behaviour.
Since promotion reduces $f$, trust was self-limiting: measured $T^\star \approx
0.29$ against $\theta_{\mathrm{high}}=0.7$, making Tier~3 unreachable at any
round count. Not being observed is not evidence of misbehaviour; it is a reason
to observe again, which \eqref{eq:stale} now expresses. The fixed point is now
$T^\star=\varphi$.}

\begin{align}
i \in \Vset:&\quad a_i \leftarrow 0,\qquad \ell_i \leftarrow \round,
\label{eq:reset}\\
i \in \Pset:&\quad a_i \leftarrow a_i + 1,
\label{eq:advance}\\[4pt]
c_i \leftarrow &\
\begin{cases}
c_i + 1, & i \in \Vset \ \wedge\ i \notin \Oset \ \wedge\ \varphi_i > 0.8,\\
0, & \text{otherwise.}
\end{cases}
\label{eq:streak}
\end{align}

The outlier condition in \eqref{eq:streak} is the detector's own verdict, not a
score cutoff: under \eqref{eq:phi} a client marginally past the anomaly
threshold still scores above $0.8$, so a flagged client would otherwise
continue accumulating the streak that Tier~3 requires.

% =============================================================================
\section{Layer 2 --- Ephemeral Structured Projection}

Per-round, per-block projection matrices are derived from a secret key, so no
client can anticipate the subspace its update will be examined in:
\begin{align}
\mathrm{seed}_m(\round) &= \mathrm{SHA\text{-}256}\bigl(\text{key} \,\|\, \round \,\|\, m\bigr),
\label{eq:seed}\\[2pt]
R_m(\round) &\sim \mathcal{N}\!\left(0, \tfrac{1}{k_m}\right)^{k_m \times d_m},
\qquad \textstyle\sum_m k_m = k,
\label{eq:projmat}\\[2pt]
\gproj_i^{(m)} &= R_m(\round)\, g_i^{(m)}.
\label{eq:project}
\end{align}

By the Johnson--Lindenstrauss lemma pairwise distances are preserved up to
$(1\pm\epsilon)$ with $k_m = O(\epsilon^{-2}\log N)$, so detection may run on
sketches rather than on full parameter blocks.

\medskip
\noindent
\textbf{Allocation caveat.} The implementation allocates $k_m$ proportionally to
$d_m$ with a floor of $1$. For the evaluated model at $k=150$ this yields
$k_m=1$ for six of ten blocks, where \eqref{eq:z} is a single squared Gaussian
and therefore extremely heavy-tailed. This is the source of the false-positive
behaviour discussed in \S\ref{sec:bvd}.

% =============================================================================
\section{Layer 3 --- Detection and Aggregation}
\label{sec:bvd}

\subsection{Block-variance outlier detection}

\begin{align}
\bar{g}^{(m)} &= \operatorname*{median}_{i \in \Vset} \gproj_i^{(m)}
\qquad\text{(coordinate-wise)},
\label{eq:center}\\[2pt]
d_i^{(m)} &= \bigl\lVert \gproj_i^{(m)} - \bar{g}^{(m)} \bigr\rVert_2^2,
\label{eq:dist}\\[2pt]
z_i^{(m)} &= \frac{d_i^{(m)}}{\sigsq_m^2 + \varepsilon},
\qquad
Z_i = \max_m z_i^{(m)},
\label{eq:z}\\[2pt]
\Oset(\round) &= \{\, i \in \Vset : Z_i > \tau_z \,\},
\qquad
\Lset(\round) = \Vset(\round) \setminus \Oset(\round),
\label{eq:outlier}\\[2pt]
\sigsq_m^2 &\leftarrow \alpha_\sigma\, \sigsq_m^2
+ (1-\alpha_\sigma)\,\frac{1}{\lvert\Lset\rvert}\sum_{i \in \Lset} d_i^{(m)}.
\label{eq:sigma}
\end{align}

\medskip
\noindent
\textbf{Known defects in \eqref{eq:z}--\eqref{eq:sigma}.} These are stated
because the empirical false-positive rate reaches $26\%$ on IID data with no
adversary present, and the cause is structural rather than a matter of tuning
$\tau_z$:
\begin{enumerate}\itemsep2pt
\item $Z_i$ is a ratio of a squared distance to a \emph{mean} squared distance,
      not a standardised score --- no dispersion estimate appears in the
      denominator, so $\tau_z$ admits no interpretation as a tail probability.
\item $\sigsq_m^2$ in \eqref{eq:sigma} is estimated only over $\Lset$, i.e. over
      clients already accepted. This truncates the upper tail of the quantity
      being estimated and biases it downward on every round.
\item $\max_m$ in \eqref{eq:z} carries no multiplicity correction, so the
      detector becomes stricter as blocks are added, and the block with the
      smallest $\sigsq_m^2$ determines every verdict. Measured divergence
      between the smallest and largest $\sigsq_m^2$ reaches $4\times10^3$ over
      $100$ rounds.
\end{enumerate}

\subsection{Behaviour score}

\begin{equation}
\varphi_i(\round)=
\begin{cases}
1, & Z_i \le \tau_z,\\[2pt]
\max\!\left(0,\ 1 - \dfrac{Z_i - \tau_z}{\tau_z}\right), & Z_i > \tau_z.
\end{cases}
\label{eq:phi}
\end{equation}

\changed{Previously $\varphi_i = 1 - d_i / \max_j d_j$, i.e.\ normalised by the
round's \emph{worst} client. One client therefore scored $0$ in every round,
including all-honest cohorts, and the cohort mean was pinned near
$1 - \bar{d}/\max_j d_j \approx 0.3$. Since \eqref{eq:trust} converges to
$\varphi$, this was a second independent mechanism capping trust below
$\theta_{\mathrm{high}}$. The score is now absolute, reusing the detector's own
anomaly threshold rather than introducing a further constant.}

\subsection{Magnitude bound on unverified updates}

\begin{align}
c &= \operatorname*{median}_{j \in \Vset} g_j
\qquad\text{(coordinate-wise)},
\label{eq:consensus}\\[2pt]
\tau &= \kappa \cdot \operatorname*{median}_{j \in \Vset} \lVert g_j - c \rVert_2,
\label{eq:radius}\\[2pt]
g_i &\leftarrow c + (g_i - c)\cdot\min\!\left(1,\ \frac{\tau}{\lVert g_i - c\rVert_2}\right),
\qquad i \in \Pset.
\label{eq:clip}
\end{align}

The radius is self-calibrating: it tracks the verified cohort's own spread, so
no absolute scale constant is required, and $\kappa=1$ states exactly that a
promoted client may deviate as much as a typical verified client does. Only
magnitude is affected --- direction is preserved and clients inside the ball are
returned unchanged. Verified clients are exempt, having already passed
\eqref{eq:outlier}.

\medskip
\noindent
\textbf{Why \eqref{eq:budget} alone is insufficient.} The budget bounds the
weight unverified clients hold, not the norm of what they submit. A measured
configuration with $\gamma = 0.114$ produced an aggregate norm of order $10^6$;
the two bounds constrain orthogonal quantities and neither subsumes the other.

\subsection{Aggregation}

\begin{align}
v_i &= \max(0.05,\ \varphi_i)\cdot n_i, \qquad i \in \Lset,
\label{eq:vweight}\\[2pt]
w_i &= p_i(\round)\cdot n_i, \qquad\qquad\ \ i \in \Pset,
\label{eq:pweight}\\[2pt]
Z(\round) &= \sum_{i \in \Lset} v_i + \sum_{i \in \Pset} w_i,
\label{eq:norm}\\[4pt]
g(\round+1) &= \frac{1}{Z(\round)}
\left( \sum_{i \in \Lset} v_i\, g_i + \sum_{i \in \Pset} w_i\, g_i \right).
\label{eq:aggregate}
\end{align}

\changed{The factor $n_i$ in \eqref{eq:vweight}--\eqref{eq:pweight} was absent:
weighting derived from trust alone. Under a Dirichlet split at $\alpha=0.3$ the
measured spread of client dataset sizes is $6.5\times$ (811 to 5261 samples), so
a client holding $811$ samples carried the same weight as one holding $5261$.
The trust term is retained as a multiplier, so unverified clients are still
discounted --- but relative to a correct FedAvg base weight rather than in place
of it. Aggregation is invariant to a common rescaling of all weights, which is
what makes this safe: clients submit full parameters, so a change in overall
scale would corrupt the global model rather than re-balance it.}

% =============================================================================
\section{Summary of divergences from the manuscript}

\begin{center}
\begin{tabular}{@{}llp{6.6cm}@{}}
\toprule
\textbf{Eq.} & \textbf{Quantity} & \textbf{Change} \\
\midrule
\eqref{eq:trust} & Trust update &
Promotion no longer decays trust; staleness counters
\eqref{eq:stale} carry that role. Fixed point moves from $f\varphi$ to
$\varphi$, making Tier~3 reachable. \\[4pt]
\eqref{eq:phi} & Behaviour score &
Absolute rather than relative to the round's worst client. \\[4pt]
\eqref{eq:vweight}--\eqref{eq:pweight} & Aggregation weights &
Dataset size $n_i$ restored (standard FedAvg). \\[4pt]
\eqref{eq:stale}, \eqref{eq:assign} & Staleness expiry &
New. Bounds the interval a promoted client may go unverified, in both
appearances and wall-clock rounds. \\
\bottomrule
\end{tabular}
\end{center}

\medskip
\noindent
Equations \eqref{eq:z}--\eqref{eq:sigma} are reported as implemented and are
known to be defective; they are the subject of ongoing revision and should not
be presented as a validated component.

\end{document}
